Os dados do \sintonia\ envolvem composição corporal e indicadores de saúde dos usuários, o
que exige cuidado com privacidade, confidencialidade e segurança das informações
armazenadas.

A participação de voluntários nas etapas de avaliação ocorre mediante consentimento livre e
esclarecido, e os dados são usados apenas para fins acadêmicos, de pesquisa e de
aperfeiçoamento da plataforma. Sempre que aplicável, as informações passam por anonimização,
que remove os elementos capazes de identificar os participantes.

O projeto observa os princípios da Resolução CNS n\textordmasculine\ 510/2016, que dispõe
sobre as normas para pesquisas em Ciências Humanas e Sociais envolvendo seres humanos, bem
como as recomendações éticas para o tratamento de informações pessoais em ambientes
digitais.

O processamento e o armazenamento dos dados seguem a Lei Geral de Proteção de Dados Pessoais
(Lei n\textordmasculine\ 13.709/2018 --- LGPD), para garantir a privacidade, a segurança e o
controle das informações pelos usuários. Como o \emph{Small Language Model} roda localmente,
há menos necessidade de compartilhar dados sensíveis com serviços externos de IA, o que
fortalece a proteção das informações de saúde.

O projeto também observa a Resolução n\textordmasculine\ 2.454/2026 do Conselho Federal de
Medicina~\cite{cfm2026}, que trata do uso de inteligência artificial na medicina: a assistente atua como
apoio, apresenta as fontes das classificações e não emite diagnóstico automático, e o
aplicativo sinaliza quando o processamento ocorreu em nuvem. Para usuários menores de 18
anos, o cadastro prevê fluxo específico com responsável legal, conforme o Art.~14 da LGPD, e
o projeto prevê a figura do encarregado pelo tratamento de dados (DPO).

As interpretações do sistema têm caráter exclusivamente informativo e educacional. O
\sintonia\ não faz diagnósticos médicos, prescrições terapêuticas nem recomendações clínicas
individualizadas, e não substitui a avaliação de profissionais de saúde habilitados.
